\Needspace{14\baselineskip}
\section{Potential regularity and metric distance}
\label{sec:psh-holder-geometry}

A H\"older modulus of continuity of the potential controls the trace mass of a positive current on small balls. Averaging the lengths of segments against this trace measure then bounds the distance of the metric defined by the current on its regular locus. This turns regularity of the potential into a metric estimate without any pointwise bound on the Hessian.

Let $U\subset\C^n$ be a coordinate domain with Euclidean form $\omega_{\mathrm E}=\ddc|z|^2$ and Euclidean volume $\mathrm{d}\lambda$. Let $0<\alpha\leq1$, let $\theta$ be a smooth real $(1,1)$-form on $U$, and let $v\in C^{0,\alpha}(U)$ satisfy $T=\theta+\ddc v\geq0$. Write $\mu=T\wedge\omega_{\mathrm E}^{n-1}$ for the trace measure of $T$.

\begin{theorem}\label{thm:holder-current-distance}\index{Holder potential@H\"older potential!trace mass}\index{Holder potential@H\"older potential!metric distance}
For every open set $V\Subset U$ there are constants $C,r_0>0$ such that
\begin{equation}\label{eq:trace-ball}
 \mu(B(x,r))\leq C r^{2n-2+\alpha},\qquad x\in V,\quad0<r<r_0.
\end{equation}
The total variation of every coefficient measure of $T$ is bounded by a constant multiple of $\mu$. Hence the coefficient measures of $T$ vanish on sets of zero Hausdorff $(2n-2+\alpha)$-measure, and all Lelong numbers of $v$ vanish. The estimate implies neither absolute continuity of $\mu$ nor a pointwise bound for the Hessian of $v$.

Suppose in addition that $T$ is a smooth Hermitian form $\eta$ on $U\setminus A$ for a proper analytic subset $A\subset U$. For $x,y\in U\setminus A$, let $d_\eta^U(x,y)$ be the infimum of the $\eta$-lengths of piecewise smooth paths in $U\setminus A$ joining $x$ and $y$. Then every point of $U$ has a neighborhood $V\Subset U$ such that, for a constant $C$,
\begin{equation}\label{eq:holder-distance}
 d_\eta^U(x,y)\leq C|x-y|^{\alpha/2},\qquad x,y\in V\setminus A.
\end{equation}
In particular, all sequences in $V\setminus A$ converging to a given point of $A\cap V$ are Cauchy for $d_\eta^U$ and have a common limit in the metric completion. If these local hypotheses hold in charts covering $X$, with a proper analytic subset $A\subset X$ and a Hermitian metric $\eta$ on $X\setminus A$, then the metric completion of $X\setminus A$ for its intrinsic $\eta$-distance is compact, and the identity of $X\setminus A$ extends to a continuous surjection of $X$ onto it. The metric $\eta$ is incomplete when $A\ne\varnothing$. Both exponents in \eqref{eq:trace-ball} and \eqref{eq:holder-distance} are sharp.
\end{theorem}

\begin{proof}
Choose a nonnegative cutoff function $\chi$ equal to one on $B(x,r)$ and supported in $B(x,2r)$, with second derivatives bounded by $Cr^{-2}$. Positivity and integration by parts, after subtracting the constant $v(x)$ from $v$, give
\[
 \mu(B(x,r))\leq C_1r^{2n}+C_2r^{2n-2}
  \operatorname{osc}_{B(x,2r)}v
 \leq Cr^{2n-2+\alpha}.
\]
Here we used that the constant $v(x)$ integrates to zero against $\ddc\chi\wedge\omega_{\mathrm E}^{n-1}$. The matrix $(a_{ij})$ of Radon--Nikodym derivatives of the coefficient measures with respect to $\mu$ is positive semidefinite almost everywhere, so $|a_{ij}|\leq(a_{ii}+a_{jj})/2$. A ball cover proves the assertion about Hausdorff null sets. After a smooth local correction making $v$ plurisubharmonic, dividing \eqref{eq:trace-ball} by $r^{2n-2}$ shows that all Lelong numbers vanish.

For the distance estimate, let $m=2n$ and let $\tau$ be the Euclidean trace of $\eta$. Then $\tau\,\mathrm{d}\lambda\leq C\mu$ on the regular locus and $|\xi|_\eta\leq C\sqrt\tau|\xi|_{\mathrm E}$ for tangent vectors $\xi$. Polar integration bounds the average $\eta$-length of the segments from $x$ to the points of $B(x,R)$ by a constant times $\int_{B(x,R)}\sqrt\tau(\zeta)|\zeta-x|^{1-m}\,\mathrm{d}\lambda(\zeta)$. On a dyadic annulus of radius $s$, the Cauchy--Schwarz inequality and \eqref{eq:trace-ball} bound the contribution to this integral by $Cs^{1-m}s^{m/2}s^{(m-2+\alpha)/2}=Cs^{\alpha/2}$. Summing over the annuli gives the bound $CR^{\alpha/2}$.

Now let $x,y\in V\setminus A$, and average the lengths of the broken segments $[x,z]\cup[z,y]$ over $z$ in the ball of radius $|x-y|/2$ about the midpoint of $x$ and $y$. By the preceding bound, some such path has length at most $C|x-y|^{\alpha/2}$. For almost every $z$, both segments avoid $A$: proper analytic sets are locally countable unions of smooth strata of real codimension at least two, and the cones from a fixed regular point over these strata have zero volume. If $V$ is small enough, all these paths stay inside $U$.

By this estimate, all regular sequences converging to a point of $A$ have a common limit in the completion, and the identity of the regular locus extends continuously across $A$. By compactness, the image of this extension is the entire completion; different points of $A$ may be identified. A sequence converging to a point of $A$ cannot converge in the metric to a regular point, since $\eta$ is bounded below by a positive multiple of the Euclidean metric on a small closed neighborhood of that point. Hence the metric $\eta$ is incomplete when $A$ is nonempty. Finally, the local model $v=|z|^\alpha+|w|^2$, in coordinates $(z,w)\in\C\times\C^{n-1}$, has trace mass of exact order $r^{2n-2+\alpha}$ on balls centered on $\{z=0\}$, and the distance to $\{z=0\}$ is of exact order $|z|^{\alpha/2}$; this proves both sharpness claims. In later applications to potentials with logarithmic poles, the theorem is applied to the positive residual current that remains after the divisorial part has been removed.
\end{proof}
